\section{Two prediction targets}\label{sec:targets}

Let $c=1,\dots,K$ index calibration populations and $j=1,\dots,d$ a fixed stack
of response thresholds, and where applicable survey rounds. Relative to a
prediction centre $\mu$,
\begin{equation}\label{eq:decomp}
Y_c(j)=\Ftilde_c(j)-\mu(j)=G_c(j)+S_c(j),\qquad G_c(j)=F_c(j)-\mu(j),
\end{equation}
with $S_c$ the survey estimation error and $v_c^2(j)$ its design variance.
Conditional on a finite population, design-unbiased or approximately unbiased
estimation motivates a centred error; for a weighted empirical distribution the
usual ratio form is not exactly design-unbiased, and we keep the idealised
centred model of the theory separate from its survey approximation.

The two statements of interest are
\begin{equation}\label{eq:t1t2}
\Pr\{\Ftilde_{K+1}(j)\in\Bhat(j)\ \forall j\}\ge 1-\alpha
\quad\text{(T1)},\qquad
\Pr\{F_{K+1}(j)\in\Bhat(j)\ \forall j\}\ge 1-\alpha
\quad\text{(T2)}.
\end{equation}
Both average over the population-level model and, for observed objects, their
sampling mechanisms. Neither is conditional coverage for a named population.
Exchangeability is assumed of the observed curves \emph{together with their
designs and observation patterns}; a shared label such as ``country'' does not
establish it, and unlike designs can defeat it even when the latent curves are
exchangeable. Departures are a separate problem with its own
framework \citep{barber2023conformal}.

Table~\ref{tab:targets} states what each construction may claim.

\begin{table}[htbp]\centering
\caption{Constructions, the target each may claim, and the basis of the claim.
All statements are simultaneous on a fixed finite grid.}\label{tab:targets}
\begin{tabular}{p{.24\linewidth}p{.20\linewidth}p{.46\linewidth}}\toprule
Construction & Target & Basis\\\midrule
Observed anchor & T1 & Exchangeable observed curves; finite-sample rank bound\\
Noise enlargement & T2 & Anchor, widened by a bound on the target's maximum sampling error\\
Oracle scale correction & T2 & Common standardised-shape law (S), with known scales\\
Estimated correction & T2 & (S), plus control of scale, centring and design-variance approximation error\\
\bottomrule\end{tabular}
\end{table}

\paragraph{Anchor.} For a centre fixed independently of the curves, let
$A_c=\max_j|Y_c(j)|$, $m=\lceil(1-\alpha)(K+1)\rceil$ and $q_\alpha=A_{(m)}$,
with $q_\alpha=\infty$ when $m>K$. Exchangeability gives T1 coverage at least
$m/(K+1)$. A radius that does not exist is reported as infinite rather than
silently replaced by the sample maximum: the band is uninformative, and saying
so is part of the guarantee.

\paragraph{Enlargement.} If a deterministic $u\ge0$ satisfies
$\Pr\{\|S_{K+1}\|_\infty>u\}\le\beta$, then a T1 band at level $1-\alpha_0$
enlarged by $u$ covers the latent curve with probability at least
$1-\alpha_0-\beta$, by the triangle inequality on the intersection of the two
events. No independence is needed and no shape is assumed. The route is
shape-free but wide, and it is not operational for an unsurveyed population
without information bounding the design error of its hypothetical survey.

\paragraph{Non-identification.} Without information restricting $S$, the law of
$G$ is not identified by that of $Y=G+S$: distinct signal and noise laws
generate the same convolution \citep{fan1991optimal}. This motivates the
additional assumptions below. It does not establish that any particular
conformal radius is optimal among latent-target procedures, and we make no
width-optimality claim.
